\section{Endocitosis}

Es un tipo de transporte activo que mueve partículas hacia una célula. Se caracteriza porque su mecanismo consiste en la invaginación de la membrana plasmática para formar un "bolsillo" alrededor de la partícula diana.

Partículas transportadas:

\begin{itemize}
    \item Moléculas grandes
    \item Partes de una célula
    \item Células enteras
\end{itemize}

Variaciones:

\begin{itemize}
    \item Fagocitosis: Proceso de "comer" células o partículas grandes.
    \item Pinocitosis: Variante donde se consumen partículas pequeñas junto con agua del líquido extracelular. No necesita de lisosomas.
    \item Endocitosis mediada por receptores
\end{itemize}

\begin{enumerate}
    \item Comienza con la proteína clatrina adhiriéndose a la membrana.
    \item Esta porción se extiende para rodear la partícula y crea una vesícula.
        \begin{itemize}
            \item Endosoma
            \item Fagosoma
        \end{itemize}
    \item Se acidifica el ambiente para "digerir" la partícula y se ponen en acción los lisosomas.
        \begin{itemize}
            \item Enzimas: Proteasas, lipasas y nucleasas
            \item Moléculas antimicrobianas
        \end{itemize}
    \item Se libera el material digerido (exocitosis).
\end{enumerate}